\begin{lemma}[Lemma 4.1, with the decay rate (4.17) of the thesis]
Let $n\ge2$, $p>1$, $0<\lambda\le\Lambda$ and $a>0$, let $0<\nu_0<\nu_1$,
suppose the Sobolev--Poincar\'e inequality on balls holds in dimension $n$
(\noderef{CitedSobolevPoincareBall}), and assume the cited Dirichlet result (C3)
with H\"older boundary data (\noderef{CitedDirichletHolder}) in the universal
form described below.  Then (A) and (B) below both hold.

\medskip
\noindent\textbf{(A) The estimate of Lemma 4.1.}  Let $M\ge0$, let $d>0$ and
let $\sigma\colon\R\to\R$.  Then there exist a nonnegative, bounded,
continuous, increasing function $\sigma_2$ with $\sigma_2(0)=0$ and a number
$R_1>0$, both depending only on
$n,p,\lambda,\Lambda,a,\nu_0,\nu_1,M,d$ and $\sigma$, such that the following
holds.  Let $0\le\varepsilon\le1$, let $\Omega\subseteq\R^n$ be open, let
$\Omega'\subseteq\R^n$ be open and bounded with
$\overline{\Omega'}\subseteq\Omega$ and $d\le d(\Omega',\partial\Omega)$ (i.e.\
$d\le|x-y|$ for every $x\in\Omega'$ and every $y\notin\Omega$).  Let $f$ satisfy
(F1) of the thesis--(F3) of the thesis on $\Omega$ with parameters
$p,\varepsilon,\lambda,\Lambda,a,\nu_0,\nu_1$ and modulus $\sigma$
(\noderef{VariationalStructureF123}), assume the cited H\"older estimates (C4)
(\noderef{CitedHolderOscillationMinimum}), and let $u\in W^{1,p}(\Omega)$ with weak
gradient $g=\nabla u$ (\noderef{MemW1pOn}), $|u|\le M$ a.e.\ on $\Omega$, be a
local minimum for $F$ in $\Omega$ (\noderef{IsLocalMinimumOn}).  Let
$x_0\in\Omega'$ and $0<R\le R_1$ with $\overline B_R(x_0)\subseteq\Omega$, and
let $v$ be the comparison map for $u$ on $B_R(x_0)$ with the frozen operator
$A=\nabla_hf^0$, $f^0(h)=f(x_0,U(x_0),h)$
(\noderef{FrozenGradientOperator}, \noderef{FrozenComparisonOn}), with weak
gradient $g_v$, continuous extension $V$ and continuous representative $U$ of
$u$ on $\overline B_R(x_0)$.  Put $w=u-v$, so $\nabla w=g-g_v$.  Then
\[
\int_{B_R(x_0)}|\nabla w|^p\,dx\;\le\;\sigma_2(R)\int_{B_R(x_0)}(1+|\nabla u|^2)^{\frac p2}dx .
\]

\medskip
\noindent\textbf{(B) The decay rate (4.17) of the thesis.}  Let $0<\alpha\le1$ and
$\delta>0$.  Then there is $\delta'>0$, depending only on
$n,p,\lambda,\Lambda,a,\nu_0,\nu_1,\alpha$ and $\delta$, with the following
property.  Let $M\ge0$, let $d>0$, let $\sigma\colon\R\to\R$, let $c_\sigma>0$
and $C_1\ge0$, and suppose the modulus obeys the H\"older condition
(F4) of the thesis,
\[
\sigma(t)\le c_\sigma t^{\delta}\qquad(t\ge0).
\]
Then there is $R_1'>0$, depending only on
$n,p,\lambda,\Lambda,a,\nu_0,\nu_1,\alpha,\delta,c_\sigma,C_1,M,d$ and
$\sigma$, such that the following holds.  Let $0\le\varepsilon\le1$, let
$\Omega\subseteq\R^n$ be open, let $\Omega'\subseteq\R^n$ be open and bounded
with $\overline{\Omega'}\subseteq\Omega$ and $d\le|x-y|$ for every
$x\in\Omega'$ and every $y\notin\Omega$.  Let $f$ satisfy
(F1) of the thesis--(F3) of the thesis on $\Omega$ with parameters
$p,\varepsilon,\lambda,\Lambda,a,\nu_0,\nu_1$ and modulus $\sigma$
(\noderef{VariationalStructureF123}), assume the cited H\"older estimates (C4)
(\noderef{CitedHolderOscillationMinimum}), and let $u\in W^{1,p}(\Omega)$ with weak
gradient $g=\nabla u$ (\noderef{MemW1pOn}), $|u|\le M$ a.e.\ on $\Omega$, be a
local minimum for $F$ in $\Omega$ (\noderef{IsLocalMinimumOn}).  Let
$x_0\in\Omega'$ and $0<R\le R_1'$ with $\overline B_R(x_0)\subseteq\Omega$, and
let $v,V,g_v,U$ be as in (A), i.e.\ let the frozen comparison bundle
$\mathrm{FrozenComparisonOn}(n,p,x_0,R,\nabla_hf^0,u,U,g,v,V,g_v)$ hold with
$f^0(h)=f(x_0,U(x_0),h)$.  Suppose in addition that the representative $U$ is
H\"older continuous of exponent $\alpha$ on the closed ball with constant
$C_1$:
\[
|U(x)-U(y)|\le C_1|x-y|^{\alpha}
\qquad\bigl(x,y\in\overline B_R(x_0)\bigr).
\]
Then
\[
\int_{B_R(x_0)}|\nabla w|^p\,dx\;\le\;R^{\delta'}\int_{B_R(x_0)}(1+|\nabla u|^2)^{\frac p2}dx .
\tag{4.17}
\]

\medskip
\noindent\textbf{Why (B) is a second clause and not a strengthening of (A), and
why $\alpha$ and $\delta$ are bound inside it.}  Lemma 4.1 as the thesis states
it at l.~1162--1167 assumes no rate for the modulus: ``there exists a bounded
continuous increasing function $\sigma_2(t)$ such that $\sigma_2(0)=0$ and
$R_1>0$ depending only on $n,p,\lambda,\Lambda,a,\nu,M$ and
$d(\Omega',\partial\Omega)$''.  Neither $\alpha$ nor $\delta$ occurs there, and
$\delta$ is the exponent of (F4) of the thesis, which Lemma 4.1 does not assume at
all.  So $\alpha$ and $\delta$ must not be in scope in (A): they are bound at
the head of (B) only, and $\sigma_2$ and $R_1$ accordingly depend on exactly the
quantities the thesis lists (our $\nu_0,\nu_1$ being its $\nu$, and $\sigma$
being named because the modulus is a parameter of this lemma rather than fixed
once and for all).  Theorem 3 (\noderef{MinimumMorreyRegularity}) consumes (A)
and nothing else; Theorem 4 consumes (B).

\medskip
\noindent\textbf{On the dependence of $\delta'$.}  The thesis records it at
l.~1359--1361: ``for a certain $\delta'>0$, depending only on
$\delta,\lambda,\Lambda,\nu_0,\nu_1,n,p$ and $a$''.  That is the list above,
with $\alpha$ adjoined, $\alpha$ being ours rather than the thesis'
because we make the H\"older exponent a bound parameter (see the next
paragraph).  In particular $\delta'$ must not depend on $M$, on $d$, on $\sigma$
or on $c_\sigma$, and it is accordingly bound before all four.  The radius
$R_1'$, by contrast, is exactly the thesis' ``radius threshold that does depend
on $c$'': it is bound after $c_\sigma$ and $C_1$ and may depend on everything.
This is what Theorem 4's own remark at l.~1353--1355 requires, where $\lambda_0$
and $\beta=\lambda_0/p$ ``depend exclusively on
$\lambda,\Lambda,\nu_0,\nu_1,n,p,a$ and $\delta$'' while the \emph{norm} depends
in addition on $d(\Omega',\partial\Omega)$, $\|u\|_{L^\infty}$ and
$\|u\|_{W^{1,p}}$.

\medskip
\noindent\textbf{Why the H\"older exponent $\alpha$ is a parameter and the
H\"older datum a hypothesis.}  The thesis derives its $\sigma_0,\sigma_1$ (and
hence the rate) from a bound $\operatorname*{osc}_{B_R}u\le2C_1R^{\alpha}$ in
which $\alpha$ is the exponent of the cited H\"older estimates (C4).  That
exponent is \emph{not} available here with the dependence $\delta'$ needs:
(C4) (\noderef{CitedHolderOscillationMinimum}) binds $M$ before its exponent, so an
$\alpha$ extracted from (C4) inside this lemma would be a function of $M$ and
so would $\delta'$.  The paper's own route avoids this, because Theorem 4 is
stated ``in the situation of Theorem 3'' and Theorem 3's part (iii)
(\noderef{MinimumMorreyRegularity}) gives $u\in C^\eta$ for \emph{every}
$\eta<1$: a consumer may therefore fix $\alpha$ --- $\alpha=\tfrac12$ serves ---
and absorb the resulting constant into $C_1$.  Accordingly $\alpha$ is bound at
the head of (B), the $C^\alpha$ datum appears only inside (B), and $C_1$ is
bound with $c_\sigma$, before $R_1'$ and after $\delta'$.  The datum is imposed
on the continuous representative $U$ and not on $u$: $u$ is a $W^{1,p}$
function, so a pointwise H\"older bound on it is not what Theorem 3 (iii)
delivers, whereas \noderef{FrozenComparisonOn} already designates $U$ as the
continuous representative and this lemma already uses the value $U(x_0)$ in the
frozen operator.

\medskip
\noindent\textbf{Two cited results, and why they are at the head.}  The proof
does \emph{not} instantiate Lemma~Lemma~3.1 of the thesis
(\noderef{DirichletProblemRestricted}).  The comparison bound it needs is
obtained in Step (6d) from \noderef{BallComparisonPrinciple}, applied to the
normalised frozen operator against constants.  That is the paper's own
disposition and not a convenience: cited result (C3) of the Note ends ``the
uniqueness and the oscillation bound of Lemma~Lemma~3.1 of the thesis (comparison
principle) must be proved'', so the comparison principle is deliberately
\emph{not} among the results this formalization may quote, and
\noderef{BallComparisonPrinciple} is the node that proves it.  The two cited
results carried at the head are therefore the following.  First, the
Sobolev--Poincar\'e inequality on balls (C6)
(\noderef{CitedSobolevPoincareBall}), because the statement of
\noderef{BallComparisonPrinciple} takes it as a hypothesis and Step (6d) must
supply it; (C6) is a statement about $n$ alone, so placing it at the head
changes no dependence list.  Second, the cited Dirichlet result (C3) with
H\"older boundary data (\noderef{CitedDirichletHolder}), carried here for
\emph{every} field $A'\colon\R^n\to\R^n$ and all parameters
$\varepsilon',\gamma_0',\gamma_1'$ with $\gamma_0'>0$ and $\gamma_1'>0$.  This
clause is the \emph{existence and regularity} half of (C3), which is the half
the Note does leave citable, and it is kept although the proof below does not
consume it: it is the hypothesis this lemma's consumers need in order to
discharge \noderef{FrozenComparisonOn}, which they do through
\noderef{FrozenComparisonDecay}, whose own statement demands exactly this
universal clause.  The universal shape is forced for a second reason as well:
the field to which (C3) would be applied is the normalised frozen operator
$\widetilde A=A^0-A^0(0)$, which is produced inside the proof at a centre and a
value chosen there, so no single instance can be named in the statement.  Both
clauses are propositions
mentioning only $n$ and $p$, both of which are already bound above them, and
each is universally quantified over its remaining parameters; so placing them at
the head leaves the dependence lists of $\sigma_2$, $R_1$, $\delta'$ and $R_1'$
exactly as stated above, and no ``depending only on'' phrase changes.  This is
the opposite of what would happen if a genuine \emph{parameter} were hoisted:
$\alpha$ and $\delta$ are real numbers that the constants may be functions of,
which is why they are bound inside clause (B) and not at the head.
\end{lemma}
\begin{proof}
Throughout, $B_\rho=B_\rho(x_0)$ and $\overline B_\rho=\overline B_\rho(x_0)$.
Fix the data of the clause being proved.  Since $x_0\in\Omega'$ and
$\overline{\Omega'}\subseteq\Omega$ we have $x_0\in\Omega$, so every clause of
\noderef{VariationalStructureF123} that is guarded by ``$x\in\Omega$'' is
available at $x=x_0$; this is used at each of (1a), (1b) and (6e) below and is
the only role of $x_0\in\Omega'$ in Steps 1--6.  Since $R>0$ we have
$x_0\in B_R$, so clause (5) of \noderef{FrozenComparisonOn} gives
$U(x_0)=u(x_0)$, and therefore
\[
f^0(h):=\mathrm{fdens}\;x_0\;(U(x_0))\;h=\mathrm{fdens}\;x_0\;(u(x_0))\;h ,
\]
which is the thesis' $f^0$ at l.~1055.  Write
$A^0=\mathrm{FrozenGradientOperator}\;n\;\mathrm{fdens}\;x_0\;(U(x_0))$
(\noderef{FrozenGradientOperator}), so that by that definition
\[
A^0(h)^i=D_{h_i}f^0(h)\qquad(h\in\R^n,\ 1\le i\le n),
\]
and $F^0(\cdot;E)=\int_Ef^0(\nabla\cdot)\,dx$.  Write
\[
D=\{h\in\R^n:\varepsilon+|h|^2\ne0\},
\]
an open set (the preimage of $\R\setminus\{0\}$ under the continuous
$h\mapsto\varepsilon+|h|^2$), and put $w=u-v$, so that $\nabla w=g-g_v$ as weak
gradients on $B_R$.

\medskip
\noindent\emph{Step 1: what the frozen integrand inherits.}

\textbf{(1a)} By clause \texttt{F1} of \noderef{VariationalStructureF123} at
$x=x_0$, for every $h$,
\[
-a+\lambda(\varepsilon+|h|^2)^{\frac p2}\le f^0(h)\le\Lambda(\varepsilon+|h|^2)^{\frac p2}+a .
\tag{1.1}
\]

\textbf{(1b)} By clause \texttt{F2diff} at $x=x_0$, $f^0$ is of class $C^1$ on
all of $\R^n$; hence $A^0$ is continuous on $\R^n$, each component being
$h\mapsto Df^0(h)[e_i]$ with $Df^0$ continuous.  By clause \texttt{F2diff2} at
$x=x_0$, $f^0$ is of class $C^2$ on $D$; so on $D$ the second derivative
$D^2f^0(h)$ exists, is continuous in $h$, and by clause \texttt{F2}, for every
$h\in D$ and every $\zeta\in\R^n$,
\[
\nu_0(\varepsilon+|h|^2)^{\frac p2-1}|\zeta|^2
\le\bigl(D^2f^0(h)\zeta,\zeta\bigr)
\le\nu_1(\varepsilon+|h|^2)^{\frac p2-1}|\zeta|^2 .
\tag{1.2}
\]
Here $(D^2f^0(h)\zeta,\zeta)$ denotes the second directional derivative
$D_\zeta\bigl(h'\mapsto Df^0(h')[\zeta]\bigr)(h)$, which is the quantity
clause \texttt{F2} bounds.

\textbf{(1c)} Put
\[
\widetilde A(h)=A^0(h)-A^0(0),
\]
the \emph{normalised} frozen field, so that $\widetilde A(0)=0$ by
construction.  By \noderef{FrozenOperatorStructureQ12} --- which is stated for
exactly this normalised field and not for $A^0$ --- the field
$(x,z,h)\mapsto\widetilde A(h)$ satisfies (Q1) of the thesis--(Q2) of the thesis on
$\R^n$ (\noderef{EllipticStructureQ12}) with the parameters
$p,\varepsilon,\gamma_0=\nu_0,\gamma_1=\nu_1$; the hypotheses of that node are
the present $2\le n$, $1<p$, $0\le\varepsilon$, $\varepsilon\le1$, $0<\lambda$,
$\lambda\le\Lambda$, $0<a$, $0<\nu_0$ and $\nu_0<\nu_1$, together with the
openness of $\Omega$, the membership $x_0\in\Omega$ recorded above, and the
\noderef{VariationalStructureF123} hypothesis of the present lemma, read at the
point $x_0$ and the value $U(x_0)$.  This is the thesis' (4.2) at l.~1073.

The bounds (Q1) of the thesis--(Q2) of the thesis constrain only
$\mathrm{fderiv}$ of the field in $h$, which subtracting the constant vector
$A^0(0)$ leaves unchanged, so (1.2) and the display above are two readings of the
same information; we use the $\widetilde A$ form wherever a node asks for
(Q1) of the thesis--(Q2) of the thesis or for $\widetilde A(0)=0$, and the $D^2f^0$ form of
(1.2) wherever the second derivative of $f^0$ itself is what appears.

\medskip
\noindent\emph{Step 2: the pointwise segment identity.}

\emph{Claim.}  For every $a_\ast,b_\ast\in\R^n$, write
$\xi(t)=t\,a_\ast+(1-t)b_\ast$ and $\zeta=a_\ast-b_\ast$, and put
\[
Q(t)=\bigl(D^2f^0(\xi(t))\zeta,\zeta\bigr)\ \text{ for }\xi(t)\in D,
\qquad Q(t)=0\ \text{ for }\xi(t)\notin D .
\]
Then $t\mapsto(1-t)Q(t)$ is interval-integrable on $[0,1]$ and
\[
f^0(a_\ast)-f^0(b_\ast)
=\bigl(A^0(b_\ast),\zeta\bigr)
+\int_0^1(1-t)\,Q(t)\,dt .
\tag{2.1}
\]

\emph{The exceptional set.}  The set
\[
T=\{t\in\R:\xi(t)\notin D\}=\{t:\varepsilon+|\xi(t)|^2=0\}
\]
contains at most one point.  Indeed $\varepsilon\ge0$ and $|\xi(t)|^2\ge0$, so
$t\in T$ exactly when $\varepsilon=0$ and $\xi(t)=0$; if $\zeta=0$ then $\xi$ is
constant and $T$ is $\R$ or $\emptyset$, and if $\zeta\ne0$ then $\xi$ is an
injective affine map and $\xi(t)=0$ has at most one solution.  In the
application made in Step 4 the two endpoints are $g(x)$ and $g_v(x)$ at a point
$x$, and the case $\zeta=0$ is the case $g(x)=g_v(x)$, in which both sides of
(2.1) vanish; so ``all but at most one $t$'' is the statement used below, and
it holds in the only case in which the integrand is not identically zero.

\emph{Proof of the claim.}  We verify the hypotheses of
\noderef{SegmentSecondOrderIdentity} and apply it with the present
$n,p,\varepsilon$, with the present $\nu_1$, with $f^0$ in the role of its $f$,
with $A^0$ in the role of its $A$, with $a_\ast,b_\ast$ in the roles of its
$a,b$, and with $Q$ in the role of its $q$.  Its hypotheses on the parameters
are $1<p$, $\varepsilon\ge0$ and $\nu_1\ge0$, all three present, the last
because $0<\nu_0<\nu_1$.

\emph{The hypotheses on $f^0$ and $A^0$.}  First, $f^0$ is differentiable at
every point of $\R^n$, being $C^1$ there by (1b).  Second, $A^0$ is continuous,
by (1b).  Third, $A^0$ represents the differential of $f^0$: by the definition
of $A^0$ (\noderef{FrozenGradientOperator}) we have
$A^0(h)^i=Df^0(h)[e_i]$, where $e_i$ is the $i$-th standard basis vector, so by
linearity of $Df^0(h)$, for every $h$ and every $\zeta'\in\R^n$,
\[
Df^0(h)[\zeta']=\sum_i\zeta'^i\,Df^0(h)[e_i]=\sum_iA^0(h)^i\zeta'^i .
\]

\emph{The hypotheses on $Q$.}  Let $t$ be such that $\xi(t)\in D$.  The set $D$
is open and $f^0$ is $C^2$ on $D$ by (1b), so $h\mapsto Df^0(h)[\zeta]$ is
differentiable on $D$ with $D_\zeta\bigl(h\mapsto Df^0(h)[\zeta]\bigr)(h)
=\bigl(D^2f^0(h)\zeta,\zeta\bigr)$ there, this being the quantity clause
\texttt{F2} bounds, as recorded in (1b).  The map $\xi$ is affine with constant
derivative $\zeta$, so by the chain rule the function
$s\mapsto\bigl(A^0(\xi(s)),\zeta\bigr)=Df^0(\xi(s))[\zeta]$ is differentiable at
$t$ with
\[
\frac{d}{ds}\Big|_{s=t}\bigl(A^0(\xi(s)),\zeta\bigr)
=\bigl(D^2f^0(\xi(t))\zeta,\zeta\bigr)=Q(t),
\tag{2.3}
\]
which is the cited node's hypothesis (i).  By the upper bound in (1.2), applied
at $h=\xi(t)\in D$ and at the vector $\zeta$,
\[
|Q(t)|\le\nu_1\bigl(\varepsilon+|\xi(t)|^2\bigr)^{\frac p2-1}|\zeta|^2 ,
\tag{2.4}
\]
which is its hypothesis (ii); here the lower bound in (1.2) is nonnegative, so
the two-sided bound there does give this bound on the absolute value.  Its
hypothesis (iii), that $Q(t)=0$ when $\xi(t)\notin D$, holds by the definition
of $Q$.

The cited node therefore gives both assertions of the claim.
\emph{End of proof of the claim.}

This is the thesis' display at l.~1089, with the two facts the thesis leaves
implicit made explicit: that the $t$-integrand is integrable even when the
segment meets the degeneracy set --- which at $\varepsilon=0$ it can, which is
why \noderef{SegmentSecondOrderIdentity} carries
\noderef{SegmentWeightIntegrable} beneath it --- and that the integration by
parts survives the single bad $t$.

\medskip
\noindent\emph{Step 3: the first-order term vanishes.}

\emph{Claim.}
\[
\int_{B_R}\bigl(A^0(g_v(x)),\,g(x)-g_v(x)\bigr)\,dx=0 .
\tag{3.1}
\]

\emph{Proof of the claim.}  This is the thesis' half-sentence at
l.~1097--1098, ``the first integral in the right hand side vanishes since $v$ is
a solution of (4.1) of the thesis and $w\in W^{1,p}_0(B_R)$'', and it is a
\emph{limit} of instances of the weak-solution property rather than an instance
of it: the test class of \noderef{IsWeakSolutionOn} requires a test pair in
$W^{1,p}(B_R)$ \emph{with compact support in $B_R$}, while clause (3) of
\noderef{FrozenComparisonOn} puts $v-u$ in $W^{1,p}_0(B_R)$
(\noderef{MemW1pZeroOn}), which supplies no compact support --- only an
approximating sequence.  That passage to the limit is
\noderef{WeakSolutionZeroTest}; what has to be supplied here is its hypothesis
on the coefficient.

\emph{The $L^{p'}$ datum.}  Let $p'=p/(p-1)$.  Apply
\noderef{FrozenFieldDualMemLp} with the present $n,p,\varepsilon$, with
$\gamma_0=\nu_0$ and $\gamma_1=\nu_1$, with $A^0$ in the role of its $A$, with
$B_R$ in the role of its $S$ and with $g_v$ in the role of its $g_v$.  Its
hypotheses are: $n\ge2$ and $p>1$, present; $0\le\varepsilon\le1$, present;
$\nu_0>0$ and $\nu_1>0$, present, the second because $0<\nu_0<\nu_1$; the
bundle
\[
\mathrm{EllipticStructureQ12}\;n\;p\;\varepsilon\;\nu_0\;\nu_1\;\R^n\;
\bigl((x,z,h)\mapsto A^0(h)-A^0(0)\bigr),
\]
which is (1c) verbatim --- the cited node is stated for the \emph{normalised}
field, which is precisely the field (1c) produces, so no transfer from
$\widetilde A$ to $A^0$ is needed here; $|B_R|<\infty$ (Preamble: the volume of
a ball); and $g_v\in L^p(B_R)$, which is clause (2) of
\noderef{FrozenComparisonOn} (\noderef{MemW1pOn}).  It gives
\[
A^0\circ g_v\in L^{p'}(B_R).
\tag{3.2}
\]

\emph{The limit.}  Apply \noderef{WeakSolutionZeroTest} with the present $n,p$,
with $B_R$ in the role of its $\Omega$, with the field
$(x,z,h)\mapsto A^0(h)$ in the role of its $A$, with $v$ and $g_v$ in the roles
of its $u$ and $g$, and with $v-u$ and $g_v-g$ in the roles of its $v$ and
$g_v$.  Its hypotheses are the following.  First $1\le n$, which follows from
$2\le n$, and $1<p$.  Second, $B_R$ is open.  Third, the weak-solution clause
\[
\mathrm{IsWeakSolutionOn}\;n\;p\;B_R\;\bigl((x,z,h)\mapsto A^0(h)\bigr)\;
\bigl((x,z,h)\mapsto0\bigr)\;v\;g_v ,
\]
which is clause (1) of \noderef{FrozenComparisonOn} verbatim: that clause
states exactly this, with the field $(x,z,h)\mapsto A^0(h)$ and the right-hand
side $0$.  Fourth, the membership of the coefficient in $L^{p'}(B_R)$: the
coefficient the node names is $x\mapsto A^0(h)$ evaluated at
$(x,v(x),g_v(x))$, and the field is constant in its first two arguments, so
this is $x\mapsto A^0(g_v(x))$, which is (3.2).  Fifth,
$\mathrm{MemW1pZeroOn}\;n\;p\;B_R\;(v-u)\;(g_v-g)$, which is clause (3) of
\noderef{FrozenComparisonOn}.  Its conclusion is
\[
\int_{B_R}\sum_iA^0(g_v(x))^i\,\bigl(g_v(x)-g(x)\bigr)^i\,dx=0 ,
\]
which is (3.1) after multiplying by $-1$.  \emph{End of proof of the claim.}

\medskip
\noindent\emph{Step 4: \textup{(4.4)} and \textup{(4.5)}.}

\textbf{(4a)} Apply Step 2 at $a_\ast=g(x)$, $b_\ast=g_v(x)$ for each $x$, so
that $\zeta=g(x)-g_v(x)=\nabla w(x)$, and integrate (2.1) over $B_R$.  Each of
the three terms is integrable on $B_R$: $f^0(g)$ and $f^0(g_v)$ by (1.1)
together with $g,g_v\in L^p(B_R)$ and $|B_R|<\infty$; the first-order term by
Step 3's computation, which exhibited $A^0\circ g_v\in L^{p'}(B_R)$ and
$\nabla w\in L^p(B_R)$; and the second-order term is then the difference of
integrable functions.  Using (3.1) to kill the first-order term,
\[
F^0(u;B_R)-F^0(v;B_R)
=\int_{B_R}\int_0^1(1-t)\bigl(D^2f^0(t\,g+(1-t)g_v)\nabla w,\nabla w\bigr)\,dt\,dx .
\tag{4.1}
\]

\textbf{(4b)} By the lower bound in (1.2), at each $x$ and for every $t$ with
$t\,g(x)+(1-t)g_v(x)\in D$ --- all but at most one $t$, by Step 2 ---
\[
(1-t)\bigl(D^2f^0(t\,g+(1-t)g_v)\nabla w,\nabla w\bigr)
\ge\nu_0(1-t)\bigl(\varepsilon+|t\,g+(1-t)g_v|^2\bigr)^{\frac p2-1}|\nabla w|^2 ,
\]
using $1-t\ge0$ on $[0,1]$.  Both sides are interval-integrable in $t$ (Step 2
for the left, \noderef{SegmentWeightIntegrable} for the right), and the
exceptional set is a single point and so null, so monotonicity of the integral
(Preamble) applies and, with (4.1),
\[
F^0(u;B_R)-F^0(v;B_R)
\ge\nu_0\int_{B_R}\int_0^1(1-t)\bigl(\varepsilon+|t\,g+(1-t)g_v|^2\bigr)^{\frac p2-1}
|\nabla w|^2\,dt\,dx ,
\tag{4.2}
\]
which is the thesis' (4.4) of the thesis.

\textbf{(4c)} Apply \noderef{MidpointGradientWeightBound} with the present $p$
--- its only hypothesis is $1<p$ --- and let $C=C(p)>0$ be the constant it
produces.  Its conclusion, applied at the present $n$ and $\varepsilon$ and at
$a=g(x)$, $b=g_v(x)$ for each $x$, with the hypotheses $1\le n$ (from $2\le n$)
and $0\le\varepsilon$, gives
\[
C\bigl(\varepsilon+|g(x)|^2+|g_v(x)|^2\bigr)^{\frac p2-1}
\le\int_0^1(1-t)\bigl(\varepsilon+|t\,g(x)+(1-t)g_v(x)|^2\bigr)^{\frac p2-1}dt .
\]
Multiplying by $|\nabla w(x)|^2\ge0$, integrating over $B_R$ and combining with
(4.2),
\[
\int_{B_R}\bigl(\varepsilon+|\nabla u|^2+|\nabla v|^2\bigr)^{\frac p2-1}
|\nabla w|^2\,dx
\le\frac1{\nu_0C}\bigl\{F^0(u;B_R)-F^0(v;B_R)\bigr\} ,
\tag{4.3}
\]
which is the thesis' (4.5) of the thesis.  This is the step the thesis calls
``elementary'' at l.~1102--1105 and whose constant it calls $C(p)$; it is
\noderef{MidpointGradientWeightBound} and nothing else is used from it.

\medskip
\noindent\emph{Step 5: minimality against the comparison map.}

\emph{Claim.}
\[
\int_{B_R}\mathrm{fdens}\;x\;(u(x))\;(g(x))\,dx
\ \le\ \int_{B_R}\mathrm{fdens}\;x\;(v(x))\;(g_v(x))\,dx .
\tag{5.1}
\]

\emph{Proof of the claim.}  As in Step 3 this is a limit, not an instance: the
test class of \noderef{IsLocalMinimumOn} is $W^{1,p}(\Omega)$ with compact
support in $\Omega$, while clause (3) of \noderef{FrozenComparisonOn} puts
$v-u$ in $W^{1,p}_0(B_R)$ (\noderef{MemW1pZeroOn}), which supplies no compact
support --- only an approximating sequence.  That passage to the limit is
\noderef{LocalMinimumZeroTest}, applied with the present
$n,p,\varepsilon,\lambda,\Lambda,a,\nu_0,\nu_1$, with $\Omega$ in the role of
its $\Omega$, with $B_R$ in the role of its $B$, with $\mathrm{fdens}$ and
$\sigma$ in the roles of its $f$ and $\sigma$, and with $u,v,g,g_v$ in the
roles of its $u,v,g,g_v$.  Its hypotheses are the following.
\begin{itemize}
\item $n\ge1$, which follows from $n\ge2$; $p>1$; $0\le\varepsilon\le1$; and
  $0<\lambda\le\Lambda$.  All four are hypotheses of the present lemma.
\item $\Omega$ is open and $B_R$ is open; $B_R\subseteq\Omega$, because
  $\overline B_R\subseteq\Omega$ is a hypothesis of the clause being proved and
  $B_R\subseteq\overline B_R$; and $|B_R|<\infty$ (Preamble: the volume of a
  ball).
\item $\mathrm{fdens}$ satisfies (F1) of the thesis--(F3) of the thesis on $\Omega$ with
  the parameters $p,\varepsilon,\lambda,\Lambda,a,\nu_0,\nu_1$ and the modulus
  $\sigma$ (\noderef{VariationalStructureF123}); this is a hypothesis of the
  clause being proved.
\item $u\in L^p(B_R)$ and $g\in L^p(B_R)$.  Both follow from
  $u\in W^{1,p}(\Omega)$ with weak gradient $g$ (\noderef{MemW1pOn}), a
  hypothesis of the clause, whose second and third clauses put $u$ and $g$ in
  $L^p(\Omega)$; restricting the measure to the subset $B_R$ preserves $L^p$
  membership.
\item $u$ is a local minimum for $F$ in $\Omega$ with weak gradient $g$
  (\noderef{IsLocalMinimumOn}); this is a hypothesis of the clause being
  proved.
\item $v-u$ lies in $W^{1,p}_0(B_R)$ with weak gradient $g_v-g$
  (\noderef{MemW1pZeroOn}); this is clause (3) of
  \noderef{FrozenComparisonOn}, which is a hypothesis of the clause being
  proved.
\end{itemize}
Its conclusion is exactly (5.1).  \emph{End of proof of the claim.}

\medskip
\noindent\emph{Step 6: from \textup{(4.5)} to \textup{(4.8)}.}

\emph{(6o) A split that keeps clause \textup{(A)}'s quantifier order honest.}
Clause (A) binds $\sigma_2$ and $R_1$ \emph{before} the inner block, and two
different things it needs live only inside that block.  The first is
$\mathrm{CitedHolderOscillationMinimum}\;n\;p$, out of which (6c-ii) below
manufactures $\alpha$ and $d_\ast$, and every one of
$\sigma_0,\sigma_1,\sigma_2,R_0,R_1$ is a function of those.  The second is
\emph{the four properties of $\sigma$}: clause (A) binds $\sigma$ as a bare
function $\R\to\R$ with no property attached, and that $\sigma$ is bounded,
continuous and increasing with $\sigma(0)=0$ is asserted only by
(F3) of the thesis, that is only by the hypothesis
$\mathrm{VariationalStructureF123}\;n\;p\;\varepsilon\;\lambda\;\Lambda\;a\;
\nu_0\;\nu_1\;\Omega\;\mathrm{fdens}\;\sigma$ inside the block.  Both are
needed before $\sigma_2$ and $R_1$ are chosen --- the moduli $\sigma_0,\sigma_1$
of (6e) are built from $\sigma$ and the threshold $R_1$ is chosen by continuity
of $\sigma_1$ at $0$ --- so both must be decided first.

Both are decidable there, because both are propositions in quantities already
bound: $\mathrm{CitedHolderOscillationMinimum}\;n\;p$
(\noderef{CitedHolderOscillationMinimum}) mentions only $n$ and $p$, all its
remaining parameters being universally quantified inside it, and the four
properties mention only $\sigma$, which clause (A) binds above $\sigma_2$ and
$R_1$.  So we decide the conjunction
\[
\mathcal P:\quad\bigl(\sigma\ \text{continuous, increasing},\ \sigma(0)=0,\
\exists C_\sigma\ \forall t\ \sigma(t)\le C_\sigma\bigr)\ \wedge\
\mathrm{CitedHolderOscillationMinimum}\;n\;p
\]
and argue by cases.

If $\mathcal P$ \emph{fails}, the inner implication of clause (A) has a false
hypothesis and is therefore satisfied by any choice; we take $\sigma_2\equiv0$
and $R_1=1$.  These meet the five demands the statement makes on them:
$R_1=1>0$; $\sigma_2$ is bounded (by $0$), continuous, nondecreasing and
nonnegative on all of $\R$; and $\sigma_2(0)=0$.  To see that the hypothesis is
indeed false, suppose the inner block's hypotheses hold.  Then
$\mathrm{CitedHolderOscillationMinimum}\;n\;p$ is one of them, and the four
properties of $\sigma$ are the four clauses \texttt{sigmacont},
\texttt{sigmamono}, \texttt{sigmazero} and \texttt{sigmabdd} of
\noderef{VariationalStructureF123}, which is another; so $\mathcal P$ holds,
contradicting the branch.  Nothing further is needed in this branch, and no
dependence list is affected, $0$ and $1$ depending on nothing at all.

If $\mathcal P$ \emph{holds}, we fix it once here and run the rest of Step 6 ---
in particular (6c-ii) and (6e) --- inside this branch.  All of
$\alpha,d_\ast,R_0,C_1,\sigma_0,\sigma_1,\sigma_2,R_1$ are then produced under
it, and their dependence lists are exactly as recorded below: neither half of
$\mathcal P$ contributes a new quantity, the first being a property of the
already-bound $\sigma$ and the second a proposition in $n$ and $p$ alone.
\emph{This is the same device the lemma block uses for
\noderef{CitedSobolevPoincareBall} and \noderef{CitedDirichletHolder}, and it is
needed here rather than there}: those two really are at the head, whereas (C4)
and (F3) are not.  Clause (B) does not need the (C4) half, because it assumes
the H\"older datum instead of deriving it, and (6c-iii) records that; but it
\emph{does} need the four properties of $\sigma$, which it binds just as bare as
clause (A) does, and the corresponding split is made there at the point where
its radius threshold is chosen.

\textbf{(6c) The H\"older datum for $U$.}  We produce $\alpha\in(0,1]$, $R_0>0$
and $C_1\ge0$, depending only on the quantities the statement allows, with
\[
|U(x)-U(x_0)|\le C_1R^{\alpha}\qquad(x\in\overline B_R,\ 0<R\le R_0).
\tag{6.3}
\]

\emph{(6c-ii) In clause \textup{(A)}.}  Apply
\noderef{MinimumHolderDatumOnBall} with the present
$n,p,\lambda,\Lambda,a,M$ and $d$.  Its hypotheses on those parameters are
$n\ge1$, which follows from $n\ge2$; $p>1$; $0<\lambda\le\Lambda$; $a>0$;
$M\ge0$; and $d>0$ --- all hypotheses of the present lemma --- together with
\noderef{CitedHolderOscillationMinimum}, which is a hypothesis of the clause
being proved and is available here because we are inside the branch of (6o) in
which it holds.  It produces $\alpha\in(0,1]$, $R_0>0$ and $C_1\ge0$ depending
only on $n,p,\lambda,\Lambda,a,M$ and $d$, and bound \emph{before}
$\varepsilon$ and before all the data; that is exactly the dependence the
statement's quantifier order demands of $\sigma_2$ and $R_1$ further down, and
it is why the cited node binds them where it does.

\emph{The subscript on the cited node's own produced constant is not
decoration.}  Part (a) of (C4) names it $d$, while $d$ in this lemma is the
domain distance $d(\Omega',\partial\Omega)$ bound in the statement; the two are
unrelated, and in particular no inequality between them is available.  The
cited node resolves this by taking our $d$ as a parameter in its own right and
writing its (C4) constant $d_\ast$; nothing of that constant escapes into the
present proof, $\alpha$, $R_0$ and $C_1$ being all that the node returns.

Its remaining hypotheses, at the data of the clause being proved, are the
following: $0\le\varepsilon\le1$; $\Omega$ open; the separation
$d\le|x-y|$ for every $x\in\Omega'$ and every $y\notin\Omega$, which is the
clause's own hypothesis; $x_0\in\Omega'$, also the clause's own;
$\mathrm{VariationalGrowthF1}\;n\;p\;\varepsilon\;\lambda\;\Lambda\;a\;
\Omega\;\mathrm{fdens}$ (\noderef{VariationalGrowthF1}), whose three clauses
\texttt{meas\_x}, \texttt{cont\_zh} and \texttt{F1} are the identically named
clauses of \noderef{VariationalStructureF123}, a hypothesis of the present
lemma; $\mathrm{IsLocalMinimumOn}\;n\;p\;\Omega\;\mathrm{fdens}\;u\;g$
(\noderef{IsLocalMinimumOn}), a hypothesis of the clause; $|u|\le M$ almost
everywhere on $\Omega$, a hypothesis of the clause; $0<R\le R_0$; that $U$ is
continuous on $\overline B_R$, which is clause (4) of
\noderef{FrozenComparisonOn}; and that $U=u$ at every point of $B_R$, which is
clause (5).  Its conclusion is
$|U(x)-U(x_0)|\le C_1R^{\alpha}$ for every $x\in\overline B_R$, which is
(6.3).

\emph{(6c-iii) In clause \textup{(B)}.}  There (6.3) is immediate from the
clause's own hypothesis $|U(x)-U(y)|\le C_1\,\mathrm{dist}(x,y)^{\alpha}$ on
$\overline B_R$, taken at $y=x_0$ and bounded by $C_1R^{\alpha}$ since
$\mathrm{dist}(x,x_0)\le R$ and $\alpha>0$; the $\alpha$ and $C_1$ are then the
clause's own, and $R_0=1$ serves.

\textbf{(6d) The comparison bound for $v$.}  We claim
\[
|v(x)-U(x_0)|\le C_1R^{\alpha}\qquad\text{for a.e.\ }x\in B_R .
\tag{6.4}
\]
This is the thesis' $|u(x_0)-v(x)|\le2C_1R^{\alpha}$ at l.~1137--1139, which the
thesis deduces from (4.3) of the thesis, $\operatorname*{osc}_{B_R}v\le
\operatorname*{osc}_{B_R}u$.  We do \emph{not} go through (4.3) of the thesis: that
statement is a clause of \noderef{DirichletProblemRestricted} about the solution
\emph{that node produces}, and identifying $v$ with it would require
$\mathrm{ContDiffOn}\;\R\;1\;v$ on $B_R$ and
$\mathrm{MemHolderOn}\;n\;\beta\;(\overline B_R)\;V$, which that node's
uniqueness clause demands and which \noderef{FrozenComparisonOn} does not
supply --- and which cannot be recovered from
\noderef{CitedDirichletHolder} either, since that node also states its
regularity only of \emph{its own} solution, so the identification would be
circular.  Instead we compare $v$ directly with constants, which is what the
bound actually needs; that comparison is \noderef{BallSolutionConstantBounds},
which also absorbs the normalisation $A^0\mapsto A^0-A^0(0)$ that the cited
comparison principle needs.

Put
\[
M_\pm=U(x_0)\pm C_1R^{\alpha}
\]
and apply \noderef{BallSolutionConstantBounds} with the present
$n,p,\varepsilon$, with $\gamma_0=\nu_0$ and $\gamma_1=\nu_1$, with $A^0$ in the
role of its $A$, with $x_0$ and $R$ in the roles of its $x_0$ and $R$, with
$v$, $V$ and $g_v$ in the roles of its $v$, $V$ and $g_v$, and with $m=M_-$ and
$M=M_+$.  Its hypotheses are the following.
\begin{itemize}
\item $n\ge2$, $p>1$, $\varepsilon\ge0$, $\nu_0>0$ and $\nu_0\le\nu_1$, the
  last from $\nu_0<\nu_1$.  All are hypotheses of the present lemma.
\item The Sobolev--Poincar\'e inequality on balls in dimension $n$
  (\noderef{CitedSobolevPoincareBall}), a hypothesis of the present lemma.
\item $R>0$, a hypothesis of the clause being proved.
\item $\mathrm{EllipticStructureQ12}\;n\;p\;\varepsilon\;\nu_0\;\nu_1\;\R^n\;
  \bigl((x,z,h)\mapsto A^0(h)-A^0(0)\bigr)$.  This is (1c) of Step 1 verbatim:
  the cited node asks for the bundle on the \emph{normalised} field, which is
  exactly the field (1c) produces, so no transfer is needed.
\item $v$ is a weak solution of $\operatorname{div}A^0(\nabla v)=0$ in $B_R$
  with weak gradient $g_v$, for the field $(x,z,h)\mapsto A^0(h)$ and the
  right-hand side $0$ (\noderef{IsWeakSolutionOn}); this is clause (1) of
  \noderef{FrozenComparisonOn}.  And $v\in W^{1,p}(B_R)$ with that same weak
  gradient (\noderef{MemW1pOn}); this is clause (2).
\item $V$ is continuous on $\overline B_R$, which is clause (6) of
  \noderef{FrozenComparisonOn}, and $V=v$ at every point of $B_R$, which is
  clause (7).
\item $M_-\le V(x)\le M_+$ for every $x\in\partial B_R$.  By clause (8) of
  \noderef{FrozenComparisonOn} we have $V=U$ on $\partial B_R$; the sphere is
  contained in $\overline B_R$, so (6.3) applies there and gives
  $|U(x)-U(x_0)|\le C_1R^{\alpha}$, that is $M_-\le U(x)\le M_+$.
\end{itemize}
Its conclusion is $M_-\le v(x)\le M_+$ for almost every $x\in B_R$, which is
(6.4).

\emph{Two remarks.}  First, (6.4) is \emph{stronger} than the thesis'
$2C_1R^{\alpha}$; we keep the thesis' constant below, since $\sigma$ is
nondecreasing and $C_1R^{\alpha}\le2C_1R^{\alpha}$, so that the coefficients of
(4.7) come out exactly as printed.  Second, (6.4) is an almost-everywhere
statement rather than a pointwise one, which is all that is used: it enters only
through the second of the two freezing errors bounded in (6e), where the
integrand is modified on a null set.

\textbf{(6e) \textup{(4.7)} and \textup{(4.8)}: the two freezing errors.}
Apply \noderef{FrozenFreezingErrorBound} with the present $n$ and $p$, with the
$\alpha$ and $C_1$ produced in (6c), with $K=\frac1{\nu_0C}$, with a bound
$C_\sigma$ for $\sigma$ and with the present $\sigma$.  Its hypotheses on those
parameters are: $n\ge1$, which follows from $n\ge2$; $p>1$;
$0<\alpha\le1$ and $C_1\ge0$, both from (6c); $K>0$, because $\nu_0>0$ and
$C>0$; and that $\sigma$ is continuous, nondecreasing, satisfies $\sigma(0)=0$
and is bounded above by $C_\sigma$ --- these four are the clauses
\texttt{sigmacont}, \texttt{sigmamono}, \texttt{sigmazero} and
\texttt{sigmabdd} of \noderef{VariationalStructureF123}, a hypothesis of the
present lemma.  It produces a constant $C_2>0$ and functions
$\sigma_0,\sigma_1\colon\R\to\R$ with
\[
\sigma_0(t)=\frac{\sigma(t+C_1t^{\alpha})+C_2\,\sigma(t+2C_1t^{\alpha})}{\nu_0C},
\qquad
\sigma_1(t)=\frac{C_2\,\sigma(t+2C_1t^{\alpha})}{\nu_0C}
\qquad(t>0),
\]
which are nonnegative, bounded above by one common constant, continuous and
nondecreasing on all of $\R$, with $\sigma_0(0)=\sigma_1(0)=0$, and which are
functions of $\alpha,C_1,K$ and $\sigma$ alone --- in particular not of
$\varepsilon$ nor of any of the data, so that they depend only on
$n,p,\lambda,\Lambda,a,\nu_0,\nu_1,M,d$ and $\sigma$, which is the dependence
the statement prescribes for $\sigma_2$ and $R_1$ below.

Its remaining hypotheses, at the data of the clause being proved, with the
present $\varepsilon$, with $z_0=U(x_0)$, with the centre $x_0$, with any
$R$ satisfying $0<R\le R_0$, with $\mathrm{fdens}$ in the role of its $f$, with
$f^0$ in the role of its $f^0$, with $u,v$ and with $g,g_v$, are the following.
\begin{itemize}
\item $g\in L^p(B_R)$ and $g_v\in L^p(B_R)$: the first because
  $u\in W^{1,p}(\Omega)$ with weak gradient $g$ (\noderef{MemW1pOn}) puts
  $g\in L^p(\Omega)$ and $B_R\subseteq\Omega$, the second by clause (2) of
  \noderef{FrozenComparisonOn}.
\item The four integrability facts, for $f^0(g)$, $f^0(g_v)$,
  $\mathrm{fdens}\;x\;(u(x))\;(g(x))$ and
  $\mathrm{fdens}\;x\;(v(x))\;(g_v(x))$ on $B_R$.  Only the first two are
  among those established in (4a); the two with $\mathrm{fdens}$ are new here
  and are supplied now.  All four share one majorant and differ only in their
  measurability.  The majorant: for every $x\in\Omega$, every $z$ and every
  $h$, clause \texttt{F1} of \noderef{VariationalStructureF123} gives
  $-a+\lambda(\varepsilon+|h|^2)^{p/2}\le\mathrm{fdens}\;x\;z\;h\le
  \Lambda(\varepsilon+|h|^2)^{p/2}+a$, and both weighted terms are nonnegative
  because $0<\lambda\le\Lambda$, so
  $|\mathrm{fdens}\;x\;z\;h|\le\Lambda(1+|h|^2)^{p/2}+a$, using
  $\varepsilon\le1$ and the monotonicity of $t\mapsto t^{p/2}$; the same bound
  holds for $f^0(h)=\mathrm{fdens}\;x_0\;(U(x_0))\;h$, which is the case
  $x=x_0$, legitimate because $x_0\in\Omega$ (Step 1).  That majorant is
  integrable on $B_R$: the constant $1$ lies in $L^p(B_R)$ because
  $|B_R|<\infty$ (Preamble: the volume of a ball), so $1+|h|\in L^p(B_R)$ for
  $h\in\{g,g_v\}$ and hence $(1+|h|)^p\in L^1(B_R)$, while
  $(1+|h|^2)^{p/2}\le(1+|h|)^p$ pointwise because $1+s^2\le(1+s)^2$ for
  $s\ge0$.  The measurability: for $f^0(g)$ and $f^0(g_v)$, the map $f^0$ is
  continuous, being $C^1$ by (1b), so the composition with an almost
  everywhere strongly measurable field is almost everywhere strongly
  measurable.  For the two with $\mathrm{fdens}$ this is not available, since
  $\mathrm{fdens}$ is only a Carath\'eodory function, and the argument is the
  following.  Clause \texttt{measx} of \noderef{VariationalStructureF123}
  makes $x\mapsto\mathrm{fdens}\;x\;z\;h$ measurable for each fixed $(z,h)$,
  and clause \texttt{contzh} makes $(z,h)\mapsto\mathrm{fdens}\;x\;z\;h$
  continuous for each fixed $x$; since $\R\times\R^n$ is a second countable
  metrisable space, the map $(z,h,x)\mapsto\mathrm{fdens}\;x\;z\;h$ is jointly
  measurable.  Replacing $u$ and $g$ by strongly measurable representatives,
  which changes $x\mapsto\mathrm{fdens}\;x\;(u(x))\;(g(x))$ only on a null set,
  and composing that jointly measurable map with the measurable
  $x\mapsto(u(x),g(x),x)$, gives the measurability of
  $x\mapsto\mathrm{fdens}\;x\;(u(x))\;(g(x))$; the same with $(v,g_v)$.  By
  domination (Preamble) all four are integrable on $B_R$.
\item The modulus bound: for every $x\in B_R$, every $z\in\R$ and every $h$,
  \[
  \bigl|\mathrm{fdens}\;x\;z\;h-f^0(h)\bigr|
  \le(1+|h|^2)^{\frac p2}\,\sigma\bigl(|x-x_0|+|z-U(x_0)|\bigr).
  \]
  This is clause \texttt{F3} of \noderef{VariationalStructureF123} at the two
  points $x$ and $x_0$ and the two values $z$ and $U(x_0)$ --- legitimate
  because $x\in B_R\subseteq\Omega$ and $x_0\in\Omega$ (Step 1) --- together
  with $f^0(h)=\mathrm{fdens}\;x_0\;(U(x_0))\;h$, which is how $f^0$ was
  defined at the head of the proof.
\item $|u(x)-U(x_0)|\le C_1R^{\alpha}$ for a.e.\ $x\in B_R$: by clause (5) of
  \noderef{FrozenComparisonOn} we have $U=u$ at every point of $B_R$, so this
  is (6.3) restricted to $B_R\subseteq\overline B_R$.  And
  $|v(x)-U(x_0)|\le C_1R^{\alpha}$ for a.e.\ $x\in B_R$: this is (6.4).
\item $\int_{B_R}\mathrm{fdens}\;x\;(u(x))\;(g(x))\,dx\le\int_{B_R}
  \mathrm{fdens}\;x\;(v(x))\;(g_v(x))\,dx$: this is (5.1) of Step 5.
\item $\int_{B_R}(\varepsilon+|\nabla u|^2+|\nabla v|^2)^{p/2-1}|\nabla w|^2dx
  \le K\bigl\{F^0(u;B_R)-F^0(v;B_R)\bigr\}$ with $K=\frac1{\nu_0C}$: this is
  (4.3) of Step 4, since $F^0(u;B_R)=\int_{B_R}f^0(g)$ and
  $F^0(v;B_R)=\int_{B_R}f^0(g_v)$ and $\nabla w=g-g_v$.
\end{itemize}
Its conclusion is
\[
\int_{B_R}\bigl(\varepsilon+|\nabla u|^2+|\nabla v|^2\bigr)^{\frac p2-1}
|\nabla w|^2dx
\le\sigma_0(R)\int_{B_R}(1+|\nabla u|^2)^{\frac p2}dx
+\sigma_1(R)\int_{B_R}|\nabla w|^p\,dx
\tag{4.8}
\]
for $0<R\le R_0$, which is the thesis' (4.8) of the thesis and the hypothesis (H)
that the absorption step below consumes.  Of the cited node's output, the
absorption step uses only the five properties of $\sigma_0$ and $\sigma_1$
listed above; the two displayed formulas are used only by clause (B), where the
same node is applied again with $\alpha_\flat$ in place of $\alpha$.

From (4.8) of the thesis the conclusion follows by the absorption step
\noderef{WeightedRemainderAbsorption}, which is where the thesis' two-regime
argument now lives, and which the thesis proves once for both lemmas.  Apply it
with the parameter $p$ --- its only hypothesis before the constants is $1<p$,
which is the present one --- and let $b_0>0$ and $c>0$ be the threshold and the
constant it produces, both depending only on $p$.  We exhibit $R_1$.  Since
$\sigma_1$ is continuous at $0$ with $\sigma_1(0)=0<b_0$, the set
$\{t:\sigma_1(t)<b_0\}$ contains a neighbourhood of $0$, hence contains some
$\eta>0$; put $R_1=\min\{R_0,\eta\}>0$, which is positive because $R_0>0$ and
$\eta>0$.  Then $\sigma_1(R_1)\le\sigma_1(\eta)\le b_0$ because $\sigma_1$ is
nondecreasing and $R_1\le\eta$.  \emph{The four properties of $\sigma$ used
here are available because we are inside the branch of \textup{(6o)} in which
$\mathcal P$ holds, which is exactly why that split was widened.}  For
$x_0\in\Omega'$ and $0<R\le R_1$,
instantiate the node with the dimension $n$, the present $\varepsilon$, with
$\sigma_0(R)$ in the role of its $a$ and $\sigma_1(R)$ in the role of its $b$,
with the ball $B_R(x_0)$ in the role of its $B$, and with $G=\nabla u=g$ and
$G_v=g_v$.  Its hypotheses hold: $0\le\varepsilon\le1$; $0\le\sigma_0(R)$ and
$0\le\sigma_1(R)\le\sigma_1(R_1)\le b_0$, the middle inequality because
$\sigma_1$ is increasing and $R\le R_1$; the ball is measurable with
$|B_R(x_0)|<\infty$; the three memberships
$|g|,|g_v|,|g-g_v|\in L^p(B_R(x_0))$ hold, the first by \noderef{MemW1pOn} on
$\Omega$ restricted to the ball, the second by the clause of
\noderef{FrozenComparisonOn} putting $v\in W^{1,p}(B_R(x_0))$ with weak
gradient $g_v$ (\noderef{MemW1pOn}), and the third by the first clause of
\noderef{MemW1pZeroOn} in the clause of \noderef{FrozenComparisonOn} that puts
$v-u\in W^{1,p}_0(B_R(x_0))$ with weak gradient $g_v-g$; and its hypothesis (H)
is (4.8) of the thesis.  Its conclusion is
\[
\int_{B_R(x_0)}|\nabla w|^p\,dx\ \le\
c\bigl(\sigma_0(R)+\sigma_0(R)^{\frac p2}+\sigma_1(R)^{\frac p2}\bigr)
\int_{B_R(x_0)}(1+|\nabla u|^2)^{\frac p2}dx ,
\]
so the lemma holds with
\[
\sigma_2(t)=c\bigl(\sigma_0(t)+\sigma_0(t)^{\frac p2}+\sigma_1(t)^{\frac p2}\bigr),
\]
and we verify its four properties, writing $C_b$ for a common bound of
$\sigma_0$ and $\sigma_1$, which the cited node of (6e) supplies, and noting
$C_b\ge\sigma_0(0)=0$.  \emph{Nonnegativity}: $\sigma_0,\sigma_1\ge0$ and
$c>0$, and a real power of a nonnegative number is nonnegative.
\emph{Boundedness}: $\sigma_2\le c(C_b+C_b^{p/2}+C_b^{p/2})$, because
$t\mapsto t^{p/2}$ is nondecreasing on $[0,\infty)$ and $p/2>0$.
\emph{Continuity}: $\sigma_0$ and $\sigma_1$ are continuous, and
$t\mapsto t^{p/2}$ is continuous on $\R$ for the exponent $p/2>0$, so each
summand is continuous and so is the constant multiple.
\emph{Monotonicity}: for $s\le t$ we get $\sigma_j(s)\le\sigma_j(t)$ and hence
$\sigma_j(s)^{p/2}\le\sigma_j(t)^{p/2}$, again because $t\mapsto t^{p/2}$ is
nondecreasing on $[0,\infty)$ and $\sigma_j\ge0$; multiplying by $c>0$
preserves the inequality.  \emph{Value at $0$}: $\sigma_0(0)=\sigma_1(0)=0$ and
$0^{p/2}=0$ since $p/2\ne0$, so $\sigma_2(0)=0$.  The
single formula the thesis displays at the end of its two-regime argument is the
$1<p<2$ branch only and is false for $p>2$; that is reported and repaired
inside \noderef{WeightedRemainderAbsorption}, which also adjudicates the
thesis' own \texttt{CHECK} on
$|\nabla w|^p\le(\varepsilon+|\nabla u|^2+|\nabla v|^2)^{p/2}$, costing a
factor $2^{p/2}$.

The quantities $\sigma_2$ and $R_1$ depend only on
$n,p,\lambda,\Lambda,a,\nu_0,\nu_1,M,d$ and $\sigma$, because $\sigma_0$,
$\sigma_1$ and $R_0$ do and because $c$ and $b_0$ depend only on $p$.

That completes clause (A).

What remains is clause (B), the decay rate (4.17) of the thesis.  \emph{Notation for
this part.}  Clause (B) binds its own H\"older exponent, which the statement
calls $\alpha$; to keep it visibly distinct from the exponent (6c-ii)
manufactured out of (C4), we write it $\alpha_\flat$ throughout the rest of the
proof, so that $0<\alpha_\flat\le1$ and the clause's datum reads
$|U(x)-U(y)|\le C_1|x-y|^{\alpha_\flat}$ on $\overline B_R$.  This matters
because the two exponents produce \emph{different} moduli: $\sigma_0,\sigma_1$
and hence $\sigma_2$ are built from $\alpha$ and $C_1$ in clause (A) and from
$\alpha_\flat$ and clause (B)'s own $C_1$ here, so they are different functions
and the thresholds chosen for them are different numbers.

\emph{(B0) The same split, for $\sigma$ alone.}  Clause (B) binds $\sigma$ as a
bare function $\R\to\R$ at the head of its block, before the threshold $R_1'$,
and its four properties are again asserted only by (F3) of the thesis inside the
inner block.  They are needed before $R_1'$ is chosen, for the same two reasons
as in clause (A): the moduli $\sigma_0^{\flat},\sigma_1^{\flat}$ are built from
$\sigma$, and the threshold is chosen by continuity of $\sigma_1^{\flat}$ at
$0$.  So we decide
\[
\mathcal P_\sigma:\quad\sigma\ \text{continuous, increasing},\ \sigma(0)=0,\
\exists C_\sigma\ \forall t\ \sigma(t)\le C_\sigma
\]
before choosing $R_1'$, which is legitimate because $\mathcal P_\sigma$
mentions only the already-bound $\sigma$.  If it fails, the inner implication
has a false hypothesis --- the four properties are the four clauses
\texttt{sigmacont}, \texttt{sigmamono}, \texttt{sigmazero} and
\texttt{sigmabdd} of \noderef{VariationalStructureF123}, which the block
assumes --- and $R_1'=1$ serves, nothing further being needed and no dependence
list being affected.  The rest of clause (B) runs inside the branch where
$\mathcal P_\sigma$ holds.  \emph{The (C4) half of (6o) is not needed here},
because clause (B) assumes the H\"older datum rather than deriving it; that is
what (6c-iii) records, and it is the only half of (6o) that clause (B) can do
without.  Note also that $\delta'$, chosen in (B3) below, needs \emph{neither}
half: it depends only on $n,p,\lambda,\Lambda,a,\nu_0,\nu_1,\alpha_\flat$ and
$\delta$, and is bound before $M$, $d$ and $\sigma$ are in scope at all, so the
split is made at the $R_1'$ choice and not at the $\delta'$ choice.

Clause (B) carries every hypothesis clause (A) carries, with the H\"older datum
adjoined, so Steps 1--5 and the derivation of
(4.8) of the thesis in (6e) all run verbatim, with (6c-iii) supplying (6.3) in place
of (6c-ii) and with $R_0=1$.  Concretely, apply
\noderef{FrozenFreezingErrorBound} again, this time with $\alpha_\flat$ in the
role of its $\alpha$ and with clause (B)'s own $C_1$, every other argument
being the one (6e) supplies; write $C_2$ for the constant it returns and
\[
\sigma_0^{\flat}(t)
=\frac{\sigma(t+C_1t^{\alpha_\flat})+C_2\,\sigma(t+2C_1t^{\alpha_\flat})}{\nu_0C},
\qquad
\sigma_1^{\flat}(t)=\frac{C_2\,\sigma(t+2C_1t^{\alpha_\flat})}{\nu_0C}
\]
for $t>0$ and $0$ for $t\le0$: these are the two functions that application
returns, by the two formulas in the cited node's conclusion, and that node also
returns their nonnegativity, their common bound, their continuity and
monotonicity on $\R$ and their vanishing at $0$.  So, for $0<R\le1$,
\[
\int_{B_R}\bigl(\varepsilon+|\nabla u|^2+|\nabla v|^2\bigr)^{\frac p2-1}
|\nabla w|^2dx
\le\sigma_0^{\flat}(R)\int_{B_R}(1+|\nabla u|^2)^{\frac p2}dx
+\sigma_1^{\flat}(R)\int_{B_R}|\nabla w|^p\,dx ,
\tag{4.8B}
\]
which is (4.8) of the thesis with the two $\flat$ moduli.

\emph{(B1) The two moduli decay at a rate the data fix.}  In (B) the H\"older
datum is a hypothesis rather than a consequence of (C4), so
$\operatorname*{osc}_{B_R}u\le C_1(2R)^{\alpha_\flat}$ directly: $U=u$ on
$B_R(x_0)$ by clause (5) of \noderef{FrozenComparisonOn}, and
$|U(x)-U(y)|\le C_1|x-y|^{\alpha_\flat}$ on $\overline B_R(x_0)$ with
$|x-y|\le2R$.  The arguments of $\sigma$ in $\sigma_0^{\flat},\sigma_1^{\flat}$
are $R+C_1R^{\alpha_\flat}$ and $R+2C_1R^{\alpha_\flat}$, exactly as the thesis
prints them for the \S5 analogue at l.~1524--1528, so the hypothesis
(F4) of the thesis of the clause --- $\sigma(t)\le c_\sigma t^{\delta}$ for $t\ge0$
--- gives, for $0<R\le1$,
\[
\sigma\bigl(R+2C_1R^{\alpha_\flat}\bigr)
\le c_\sigma\bigl(R+2C_1R^{\alpha_\flat}\bigr)^{\delta}
\le c_\sigma2^{\delta}\bigl(R^{\delta}
+(2C_1)^{\delta}R^{\alpha_\flat\delta}\bigr)
\le K_0R^{\alpha_\flat\delta} ,
\]
using $(x+y)^{\delta}\le2^{\delta}(x^{\delta}+y^{\delta})$ for $x,y\ge0$, then
$\alpha_\flat\le1$ with $R\le1$ to replace $R^{\delta}$ by
$R^{\alpha_\flat\delta}$, and the same bound with $C_1$ for $2C_1$ a fortiori;
so $K_0=c_\sigma2^{\delta}\bigl(1+(2C_1)^{\delta}\bigr)$ serves.  Dividing by
$\nu_0C$ and adding the two terms of $\sigma_0^{\flat}$, with
$K_1=(1+C_2)K_0/(\nu_0C)$,
\[
\sigma_0^{\flat}(R)\le K_1R^{\alpha_\flat\min\{\delta,1\}},
\qquad
\sigma_1^{\flat}(R)\le K_1R^{\alpha_\flat\delta}
\qquad(0<R\le1) ;
\]
the second is the display above divided by $\nu_0C$, and the first is that same
bound weakened by $R^{\alpha_\flat\delta}\le R^{\alpha_\flat\min\{\delta,1\}}$,
which holds because $\min\{\delta,1\}\le\delta$ and $R\le1$.  Here $K_0,K_1$ are
built from
$n,p,\lambda,\Lambda,a,\nu_0,\nu_1,c_\sigma,C_1,M,d$ and $\sigma$.  Note the
exponents involve only $\alpha_\flat$ and $\delta$; every dependence on $M$, $d$
and $c_\sigma$ sits in $K_1$.

\emph{(B2) The absorption, re-run in this branch.}  Apply
\noderef{WeightedRemainderAbsorption} with the parameter $p$ --- its only
hypothesis before the constants is $1<p$ --- and let $b_0>0$ and $c>0$ be the
threshold and the constant it produces, both depending only on $p$.  These are
the same two numbers clause (A) used, since they depend on $p$ alone; the
\emph{threshold radius} is not, and must be chosen for
$\sigma_1^{\flat}$ rather than for clause (A)'s $\sigma_1$.  Choose
\[
R_1''\in(0,1]\quad\text{with}\quad\sigma_1^{\flat}(R_1'')\le b_0 ,
\]
which is possible because $\sigma_1^{\flat}$ is continuous and nondecreasing on
$[0,\infty)$ with $\sigma_1^{\flat}(0)=0$ --- all three supplied by the cited
node of (6e), and available here because we are inside the branch of (B0) in
which $\mathcal P_\sigma$ holds: continuity at $0$ gives an $\eta>0$
with $\sigma_1^{\flat}(t)\le b_0$ for $0\le t\le\eta$, and $R_1''=\min\{\eta,1\}$
serves, positive because $\eta>0$.  For $0<R\le R_1''$ instantiate the node exactly as in clause (A) ---
the dimension $n$, the present $\varepsilon$, $\sigma_0^{\flat}(R)$ in the role
of its $a$ and $\sigma_1^{\flat}(R)$ in the role of its $b$, the ball
$B_R(x_0)$ in the role of its $B$, $G=\nabla u=g$ and $G_v=g_v$ --- the
hypotheses being discharged by the same citations as there, since clause (B)
carries the same data: $0\le\varepsilon\le1$; $0\le\sigma_0^{\flat}(R)$ and
$0\le\sigma_1^{\flat}(R)\le\sigma_1^{\flat}(R_1'')\le b_0$, the middle
inequality because $\sigma_1^{\flat}$ is nondecreasing and $R\le R_1''$; the
ball measurable with $|B_R(x_0)|<\infty$; the three memberships
$|g|,|g_v|,|g-g_v|\in L^p(B_R(x_0))$ by \noderef{MemW1pOn} and by the clauses of
\noderef{FrozenComparisonOn} and \noderef{MemW1pZeroOn} cited in clause (A);
and its hypothesis (H), which is (4.8B).  Its conclusion is
\[
\int_{B_R(x_0)}|\nabla w|^p\,dx\ \le\ \sigma_2^{\flat}(R)
\int_{B_R(x_0)}(1+|\nabla u|^2)^{\frac p2}dx
\qquad(0<R\le R_1''),
\tag{4.9B}
\]
with
\[
\sigma_2^{\flat}(t)=c\bigl(\sigma_0^{\flat}(t)+\sigma_0^{\flat}(t)^{\frac p2}
+\sigma_1^{\flat}(t)^{\frac p2}\bigr).
\]

\emph{The rate survives the absorption.}  Put
\[
\gamma=\alpha_\flat\min\{\delta,1\}\min\Bigl\{1,\frac p2\Bigr\}>0 ,
\]
which depends only on $p,\alpha_\flat,\delta$.  For $0<R\le1$ each of the three
terms of $\sigma_2^{\flat}(R)$ is at most a constant times $R^{\gamma}$, by
(B1): the exponents occurring are $\alpha_\flat\min\{\delta,1\}$,
$\alpha_\flat\min\{\delta,1\}\frac p2$ and $\alpha_\flat\delta\frac p2$, each of
which is at least $\gamma$, and $R\le1$ makes $R^{s}$ nonincreasing in $s$.
Hence
\[
\sigma_2^{\flat}(R)\le KR^{\gamma}\qquad(0<R\le1),
\qquad K=K(n,p,\lambda,\Lambda,a,\nu_0,\nu_1,c_\sigma,C_1,M,d,\sigma)>0 .
\]

\emph{(B3) The threshold absorbs $K$, which is what frees $\delta'$.}  Put
\[
\delta'=\frac\gamma2>0,
\qquad
R_1'=\min\bigl\{R_1'',\ 1,\ K^{-2/\gamma}\bigr\}>0 .
\]
For $0<R\le R_1'$ one has $KR^{\gamma}\le R^{\gamma/2}$, since
$R^{\gamma/2}\le(K^{-2/\gamma})^{\gamma/2}=K^{-1}$.  Since $R_1'\le R_1''$, the
radius $R$ is in the range of (4.9B), and combining the two,
\[
\int_{B_R}|\nabla w|^p\,dx\le\sigma_2^{\flat}(R)
\int_{B_R}(1+|\nabla u|^2)^{\frac p2}dx
\le R^{\delta'}\int_{B_R}(1+|\nabla u|^2)^{\frac p2}dx ,
\]
which is (4.17) of the thesis.  Note that this cites (4.9B), the conclusion of
the absorption \emph{in this branch}, and not clause (A)'s conclusion: clause
(A)'s carries $\sigma_2$, built from $\alpha$, which (B1) and (B2) do not
estimate and which would reintroduce the dependence on $M$ through (C4) that
clause (B) exists to avoid.  The exponent $\delta'=\gamma/2$ depends only on
$p,\alpha_\flat$ and $\delta$, and the whole dependence on $c_\sigma$, $C_1$,
$M$, $d$ and $\sigma$ has been pushed into the threshold $R_1'$ --- which is the
thesis' ``a radius threshold that does depend on $c$''.
\end{proof}
